\scanpage{59}
\mspage{15}
to \deleted{translate} translate the topological Teichmüller groupoid, one will never find anything but structures in which the cutting circles are actually marked.
\deleted{If one henceforth denotes by $S\mathcal TD_{\mathrm{top}}^\pm$ the groupoid}
\inserted{Nevertheless one} wants to have the right to ``erase'' them, to forget them \inserted{\unsure{these|the} markings} --- \deleted{even} if they are indeed there!

D) \uline{Gluing operation}. It is the most \unsure{crucial|essential}, the least tautological perhaps of all the operations we are reviewing.
To perform it, one needs as a preliminary datum

a) A subset \deleted{\unsure{the}} $A'$ of $A_l$ (free edges), and an involution without fixed points (that is, an equivalence relation with classes of cardinality 2) $\sigma'$ of \deleted{$\unsure{\{c,c'\}}$} $A'$.
\deleted{\unsure{and the}\ $\unsure{\sigma}$}

b) For \deleted{$\{a,a'\}\in\unsure{A'_l|A'_2}$} \unsure{every|every} orbit \deleted{of} $\{a,a'\}$ of $A'$, an orientation-reversing isomorphism from $C_a$ onto $C_{a'}$ and from $C_{a'}$ onto $C_a$, inverse to one another.
For the operation to be well determined up to unique isomorphism
\mspage{58}
